% theory/revision/thm12iii.tex -- Task 4 (G7-8): Theorem 1.2(iii) with the diameter dependence
% of the constant explicit, and "attained" qualified.  Verified by check_thm12iii.py.
%
% The constant comes from Theorem 4.9(b): on 0 < t <= l^2/(2(1+l)) the factor 1 + 2t/(l-t) is
% increasing in t, with maximum (2+3l)/(2+l) < 3 at the endpoint, and pi/sqrt(4 pi) = sqrt(pi)/2.
% Notation: the diameter is written diam, not delta (G7-31: delta is also the data error).

% ---- (1) Replace Theorem 1.2(iii) (manuscript l. 140) by:

\item Let $\Orb_1,\Orb_2$ have the same signature and area $A$, let $\ell$ be the smaller of their
systoles and $D$ the larger of their diameters. For $0<t\le\ell^2/(2(1+\ell))$,
\[
  |\cZ_{\Orb_1}(t)-\cZ_{\Orb_2}(t)|\le C(A,\ell,D)\,t^{-1/2}e^{-\ell^2/4t},\qquad
  C(A,\ell,D)=\frac{\sqrt\pi\,e^{3D}\,\ell\,e^{\ell/2}\,(2+3\ell)}{2A\,(1-e^{-\ell})\,(2+\ell)} .
\]
The constant depends on the diameter, which the signature does not determine. If the two length
spectra, counted with the weights $\ell(\gamma_0)$, first differ at the length $\ell$ (in particular,
if the two systoles differ), then $\sqrt t\,e^{\ell^2/4t}|\cZ_{\Orb_1}(t)-\cZ_{\Orb_2}(t)|$ has a
nonzero limit as $t\downarrow0$: the exponent $\ell^2/4$ and the factor $t^{-1/2}$ are attained,
the constant $C$ is not claimed to be. If they first differ at a length $L_*>\ell$, the difference
is of the smaller order $t^{-1/2}e^{-L_*^2/4t}$; if the weighted length spectra never differ
(equivalently, the orbifolds are isospectral) the difference vanishes identically. Here the
systole is the least length of a closed geodesic, including those through cone points.

% ---- (2) The sentence after Theorem 1.2 (manuscript l. 144) is unchanged; add after it:
%   "The diameter enters only through the counting bound of Lemma~\ref{lem:counting}."

% ---- (3) Theorem 4.9(b): state the same simplified bound, keeping the proof.  Replace the
%      display of (b) by
%
%   For $0<t\le\ell^2/(2(1+\ell))$,
%   \[
%     |\cZ_{\Orb_1}(t)-\cZ_{\Orb_2}(t)|\le\frac{\pi e^{3D}}{A(1-e^{-\ell})}\;\ell e^{\ell/2}\Big(1+\frac{2t}{\ell-t}\Big)\frac{e^{-\ell^2/4t}}{\sqrt{4\pi t}}
%     \le C(A,\ell,D)\,t^{-1/2}e^{-\ell^2/4t},
%   \]
%   with $D=\max(\operatorname{diam}\Orb_1,\operatorname{diam}\Orb_2)$ and $C$ as in Theorem 1.2(iii);
%
%   and add one sentence at the end of the proof of (b):
%   "On this range $t<\ell/2$, and $1+2t/(\ell-t)$ increases in $t$ to $(2+3\ell)/(2+\ell)$ at
%    $t=\ell^2/(2(1+\ell))$; with $\pi/\sqrt{4\pi}=\sqrt\pi/2$ this gives the second inequality."
%
%   Rename delta -> D (diameter) throughout Section 4 (Lemma 4.8, Theorem 4.9, Lemma 4.7 does not
%   use it).

% ---- (4) Abstract: replace "and this rate is attained." by
%   "and this rate is attained whenever the two length spectra first differ at the shorter
%    systole, for instance when the systoles differ."

% ---- (5) Section 1 pointer sentence (manuscript l. 144): "These are Theorem 4.1, Corollary 4.5
%   and Theorems 4.9 and 4.10, with Corollary 4.11." stays; the attainment clause is
%   Theorem 4.10(ii)--(iii) and Corollary 4.11, whose hypothesis L_* = ell is exactly the
%   condition stated above.

% ---- OPTIONAL (not done here): a counting bound depending on (area, systole) only would remove
%   D from C.  Not attempted; G7-8 lists it as optional.
